\section{From windows to requirements: robustness, Monte Carlo and the shooter specification}
\label{sec:robust_spec}

% Source of truth: src/shooter_spec.py (grid, variants, best_square, fixed_elev, describe,
% spec_rows, longest_run, tunable_rows, wheel_rpm, spin_rpm), src/robustness.py (wilson,
% integrate3d, backspin_axis, mc_shots, mc_scenarios, smooth_prob, best_over_hw,
% lateral_factor, wind_drift_sd, position_prob, clean_masks) and src/feasibility.py
% (integrate_batch, evaluate_gate, grid_from_candidates).
% Every worked number below was computed on 2026-10-30 with the system python against the
% repo code. Spec-grid numbers use research/params_lit.json (shooter_spec.py always loads it).
% The Monte Carlo numbers in data/robust_mc.csv and data/robust_map_*.csv were produced with the
% older data/aero_params.json (their CSV headers say so); see the comments in
% Section~\ref{sec:robust_spec-mc}.

The earlier sections give, for one release state, a yes-or-no answer: does the ball enter the
upward-facing CELL of the HIVE cleanly or not?
A shooter, though, is not a single release state.
Its exit speed and its angle scatter from shot to shot, the ball's mass and its aerodynamic
coefficients are only known within assumed, literature-based ranges, and the HIVE itself sits on the field
with a tolerance.
This section turns the yes-or-no windows into numbers a CAD designer can build to.
It goes in five steps:
(i) sample the clean window on a grid;
(ii) find the largest clean square in that grid and read it as a $\pm$ tolerance;
(iii) ask which hood angle stays usable when the aerodynamics change and the team re-tunes only
the wheel speed;
(iv) smooth the window with Gaussian release errors, and run a full three-dimensional Monte Carlo
with per-shot physics, wind and HIVE pose, scored with a Wilson interval;
(v) convert the result into wheel and ball rpm and a CAD card.

All targets at the end of this section are model outputs, so they carry \predicted.
They are design targets to be checked on the field, not measured performance.

Throughout, the frame is the one of Section~\ref{sec:flight}: $x$ is horizontal toward the HIVE,
$y$ is horizontal and sideways, and $z$ points up.
The shot comes from the unfavourable side: the ball has to pass over the roof of the HIVE and then
drop into the CELL, whose mouth faces outward and upward at the tilt angle
$\beta = 30^\circ$\confirmed\ (Competition Manual figure and official CAD).

% ---------------------------------------------------------------------------------------------
\subsection{The clean window on a grid}
\label{sec:robust_spec-grid}

\paragraph{Motivation.}
The flight solver answers ``clean or not'' for one release state.
To talk about tolerances we need that answer on a regular grid of release states, so that
neighbouring cells can be compared.

\paragraph{Release state.}
A shot is released from horizontal distance $d$ (\unit{m}) before the HIVE pivot plane and at
height $h$ (\unit{m}) above the tiles, with exit speed $v$ (\unit{m/s}), elevation angle
$\theta$ (\unit{deg}, above horizontal) and spin parameter $S$ (dimensionless, $S=\omega R/|\vect{v}|$
with ball radius $R$ in \unit{m} and spin rate $\omega$ in \unit{rad/s}).
In two dimensions there is no wind, so the forces depend only on the velocity and the spin, not on
the position.
The solver \code{integrate\_batch} therefore integrates each path once from the origin, giving
$\vect{r}_0(t;v,\theta,S)$, and translates it:
\begin{equation}
  \vect{r}(t;h,d,v,\theta,S) = \vect{r}_0(t;v,\theta,S) + (-d,\; h)^{\mathsf T}.
  \label{eq:robust_spec-translate}
\end{equation}
This is exact because the right-hand side of the equations of motion contains no position
(it depends on the velocity and, through the spin decay, on the time since release only).
One integration serves every $(h,d)$ pair, which is what makes a dense grid affordable.

\paragraph{Clean indicator.}
Let $\vect{n}_{\mathrm{out}}$ be the outward unit normal of the CELL mouth (dimensionless) and let
$c(t)$ (\unit{m}) be the raster clearance between the ball surface and the nearest HIVE part.
Following \code{evaluate\_gate}, a path is clean when it crosses the mouth edge (the gate segment
between the two lips) at a time $t_g$ (\unit{s}), moving inward, without touching anything on the
way in:
\begin{equation}
  M(h,d,v,\theta) =
  \begin{cases}
    1, & \text{if } \exists\, t_g:\ \vect{r}(t_g)\in\text{gate},\ \
         \vect{v}(t_g)\cdot\vect{n}_{\mathrm{out}} < 0,\ \ c_{\min}(t_g) \ge 0, \\
    0, & \text{otherwise,}
  \end{cases}
  \label{eq:robust_spec-mask}
\end{equation}
where $c_{\min}(t_g)=\min_{t_{\mathrm{box}}\le t\le t_g} c(t)$ (\unit{m}) and $t_{\mathrm{box}}$ (\unit{s})
is the time the ball enters the raster box around the HIVE.

\paragraph{Grid.}
The specification grid of \code{shooter\_spec.grid} is
\begin{equation}
  v_i = v_{\min} + i\,\Delta v,\qquad \theta_j = \theta_{\min} + j\,\Delta\theta ,
  \label{eq:robust_spec-grid}
\end{equation}
with $v_{\min}=5.0$\,\unit{m/s}\assumed, $\Delta v = 0.05$\,\unit{m/s}\assumed\ and 111\derived\ speeds up
to 10.5\,\unit{m/s}\assumed; $\theta_{\min}=62^\circ$\assumed, $\Delta\theta=0.1^\circ$\assumed\ and
261\derived\ elevations up to $88^\circ$\assumed.
The release heights are 0.20, 0.35, 0.50, 0.65 and 0.737\,\unit{m}\assumed, and distances run from
0.60\,\unit{m}\assumed\ to 1.75\,\unit{m}\assumed\ in steps of 0.05\,\unit{m}\assumed\ (24\derived\ values).
The spin values are $S\in\{0, 0.5, 1.0, 1.5, -1.0\}$\assumed, where negative $S$ is topspin.

The flight solver returns candidate distances $d_k$ (\unit{m}) at which a path hits the gate.
\code{grid\_from\_candidates} marks a grid distance $d$ clean only if two candidates at consecutive
time samples $k$, $k+1$ of the same path are both clean and bracket it:
\begin{equation}
  M(d) = 1 \iff \exists k:\ M(d_k)=M(d_{k+1})=1,\ \ \min(d_k,d_{k+1}) \le d \le \max(d_k,d_{k+1}).
  \label{eq:robust_spec-dcell}
\end{equation}
This is conservative: a clean point that is isolated between two dirty neighbours is not counted.

\paragraph{Worked example.}
Table~\ref{tab:robust_spec-nclean} lists the number of clean $(v,\theta)$ cells
$N(h,d)=\sum_{i,j}M(h,d,v_i,\theta_j)$ for the nominal aerodynamic set of
\code{research/params\_lit.json}.
At $d\le0.70$\,\unit{m}\assumed\ no cell is clean at the heights
$h=0.50$ to $0.737$\,\unit{m}\assumed: the ball cannot both clear the roof and descend steeply enough
into the mouth from that close.
The count peaks near $d=0.80$\,\unit{m}\predicted\ and then falls slowly with distance.
At $d=0.75$\,\unit{m}\assumed\ the window shrinks with height: 378\predicted\ cells at
$h=0.737$\,\unit{m}\assumed, 286\predicted\ at 0.65\,\unit{m}\assumed\ and 169\predicted\ at
0.50\,\unit{m}\assumed\ (no spin).
For $S=1$ at $h=0.50$\,\unit{m}\assumed\ the cell count at $d=0.75$\,\unit{m}\assumed\ is already
0\predicted.

\begin{table}[htbp]
  \centering
  \caption{Clean cells on the $111\times261$ speed--elevation grid at $h=0.737$\,\unit{m}, nominal
    aerodynamics (\code{research/params\_lit.json}).}
  \label{tab:robust_spec-nclean}
  \begin{tabular}{@{}lrrrrrrr@{}}
    \toprule
    $d$ (\unit{m}) & 0.75 & 0.80 & 0.85 & 0.90 & 1.15 & 1.50 & 1.75 \\
    \midrule
    $N$, $S=0$ & 378\predicted & 457\predicted & 436\predicted & 426\predicted & 359\predicted & 287\predicted & 247\predicted \\
    $N$, $S=1$ & 292\predicted & 412\predicted & 408\predicted & 386\predicted & 335\predicted & 244\predicted & 209\predicted \\
    \bottomrule
  \end{tabular}
\end{table}

Backspin does not make the window larger here: $S=1$ has fewer clean cells than $S=0$ at every
distance in the table.
The case $S=1.5$ gives exactly the same window as $S=1.0$, because the lift coefficient has
already saturated (Section~\ref{sec:flight}): with the nominal set, $C_L = \min(c_1 S, C_{L,\max})$ stops
growing at $S_{\mathrm{sat}} = C_{L,\max}/c_1 = 0.25/0.35 = 0.714$\derived.
More spin above that adds nothing.

\paragraph{Assumptions.}
\begin{itemize}
  \item The 2D model has no wind and no sideways motion; azimuth is added in
    Section~\ref{sec:robust_spec-azimuth}.
  \item A grid cell is a point sample; nothing is claimed about the space between grid points.
  \item The aerodynamic coefficients are assumed ranges based on pickleball and smooth-sphere data,
    not proven bounds and not POLLEN measurements.
  \item The HIVE tilt is the nominal $30^\circ$\confirmed\ unless a variant says otherwise.
\end{itemize}

% ---------------------------------------------------------------------------------------------
\subsection{The largest clean square and the $\pm$ tolerance}
\label{sec:robust_spec-square}

\paragraph{Motivation.}
A shooter set to $(v_c,\theta_c)$ that actually fires anywhere in
$[v_c-\delta_v, v_c+\delta_v]\times[\theta_c-\delta_\theta,\theta_c+\delta_\theta]$ is clean every time
only if that whole box is clean.
The largest all-clean $k\times k$ block of cells is therefore the best repeatability target the
window allows.
\code{best\_square} finds it in $O(N_v N_\theta)$ work per size using a prefix sum.

\paragraph{Prefix sum.}
Let $M_{ij}\in\{0,1\}$ be the mask on the $N_v\times N_\theta$ grid for fixed $(h,d,S)$.
Define the two-dimensional prefix sum (a count, dimensionless)
\begin{equation}
  P_{ij} = \sum_{i'=0}^{i-1}\sum_{j'=0}^{j-1} M_{i'j'},\qquad P_{0j}=P_{i0}=0 .
  \label{eq:robust_spec-prefix}
\end{equation}

\paragraph{Block sum, step by step.}
The number of clean cells in the $k\times k$ block whose lower corner is $(i_0,j_0)$ is
$B_k(i_0,j_0)=\sum_{i=i_0}^{i_0+k-1}\sum_{j=j_0}^{j_0+k-1}M_{ij}$.
The sum $P_{i_0+k,\,j_0+k}$ covers the rectangle $[0,i_0+k)\times[0,j_0+k)$.
Subtracting $P_{i_0,\,j_0+k}$ removes the rows below $i_0$, and subtracting $P_{i_0+k,\,j_0}$ removes the
columns left of $j_0$.
The corner $[0,i_0)\times[0,j_0)$ has now been removed twice, so it is added back once
(inclusion--exclusion):
\begin{equation}
  B_k(i_0,j_0) = P_{i_0+k,\,j_0+k} - P_{i_0,\,j_0+k} - P_{i_0+k,\,j_0} + P_{i_0,\,j_0}.
  \label{eq:robust_spec-block}
\end{equation}
Because each $M_{ij}\le1$, the block is all clean exactly when its count is the largest possible:
\begin{equation}
  \text{block clean} \iff B_k(i_0,j_0) = k^2 .
  \label{eq:robust_spec-allclean}
\end{equation}

\paragraph{Monotone search.}
Every $k\times k$ clean block contains a $(k-1)\times(k-1)$ clean block, so if no block of size $k$
exists, none of a larger size exists either.
The code increases $k$ from 1 and stops at the first size with no hit:
\begin{equation}
  k^\ast = \max\{k:\ \exists (i_0,j_0),\ B_k(i_0,j_0)=k^2\}.
  \label{eq:robust_spec-kstar}
\end{equation}
When several blocks of size $k^\ast$ exist, the code keeps the one whose lower corner is closest (in
index units) to the centroid $(\bar\imath,\bar\jmath)$ of the lower corners of all hits (the same as
comparing block centres, since every centre is the corner shifted by $(k^\ast-1)/2$), which keeps the
choice away from the edge of a cluster:
\begin{equation}
  (i_0,j_0) = \arg\min_{\text{hits}} \Big[\big(i_0-\bar\imath\big)^2
              + \big(j_0-\bar\jmath\big)^2\Big].
  \label{eq:robust_spec-tie}
\end{equation}

\paragraph{Centre and tolerance.}
The $k$ cell centres of the block span $k-1$ grid steps.
The set point is the middle of that span, and the half-span is the tolerance:
\begin{align}
  v_c &= v_{\min} + \Big(i_0+\frac{k^\ast-1}{2}\Big)\Delta v, &
  \theta_c &= \theta_{\min} + \Big(j_0+\frac{k^\ast-1}{2}\Big)\Delta\theta,
  \label{eq:robust_spec-centre}\\
  \delta_v &= \frac{k^\ast-1}{2}\,\Delta v, &
  \delta_\theta &= \frac{k^\ast-1}{2}\,\Delta\theta .
  \label{eq:robust_spec-tol}
\end{align}
Here $v_c$ and $\delta_v$ are in \unit{m/s}, $\theta_c$ and $\delta_\theta$ in degrees.
A single clean cell ($k^\ast=1$) gives a zero tolerance: the set point works only if the shooter is
perfect.

\paragraph{Worked example.}
At $h=0.737$\,\unit{m}\assumed, $d=0.75$\,\unit{m}\assumed\ and $S=0$\assumed, nominal aerodynamics,
the search ends at $k^\ast=3$\predicted\ with $(i_0,j_0)=(70,218)$\predicted.
Then
$v_c = 5.0 + (70+1)(0.05) = 8.55$\,\unit{m/s}\predicted\ and
$\theta_c = 62 + (218+1)(0.1) = 83.9^\circ$\predicted, with
$\delta_v = 1\times0.05 = 0.05$\,\unit{m/s}\derived\ and
$\delta_\theta = 1\times0.1 = 0.10^\circ$\derived.
The smallest raster clearance inside the block is 4.9\,\unit{mm}\predicted.
For $S=0.5$\assumed\ the square is $k^\ast=4$\predicted\ at (9.125\,\unit{m/s}\predicted,
$82.35^\circ$\predicted), above the $8.909$\,\unit{m/s}\predicted\ grip limit of
Section~\ref{sec:launcher:grip} (the search applies no speed cap); for $S=1.0$\assumed\ it is
$k^\ast=4$\predicted\ at (8.875\,\unit{m/s}\predicted, $81.85^\circ$\predicted).
Over every unfavourable-side row of the nominal sweep (nominal aerodynamics, $30^\circ$\confirmed\ tilt)
with $S\in\{0,0.5,1\}$, the square sizes are
$k^\ast=4$ in 14\predicted\ rows, $k^\ast=3$ in 250\predicted, $k^\ast=2$ in 42\predicted,
$k^\ast=1$ in 2\predicted\ and empty in 52\predicted.
No row of it reaches $k^\ast=5$.
So, for the nominal aerodynamics and tilt, $\pm0.075$\,\unit{m/s}\predicted\ and $\pm0.15^\circ$\predicted\
is the widest tolerance the grid admits (a requirement on the shooter, not a measured repeatability).
Only backspin reaches it ($S=0.5$ in 4\predicted\ rows, $S=1.0$ in 10\predicted); the no-spin twin stops
at $k^\ast=3$\predicted, i.e.\ $\pm0.05$\,\unit{m/s}\predicted\ and $\pm0.10^\circ$\predicted.
Single changed conditions can do better: with high drag at the $31.77^\circ$\derived\ tilt
($h=0.737$\,\unit{m}\assumed, $d=0.75$\,\unit{m}\assumed) backspin reaches $k^\ast=5$\predicted\ at
9.25\,\unit{m/s}\predicted, above the grip limit, and the twin $k^\ast=4$\predicted\ at
(8.875\,\unit{m/s}\predicted, $83.35^\circ$\predicted).
Each square is found for one set of coefficients and one tilt; no single setting is clean under
all conditions at once (Section~\ref{sec:robust_spec-hood}).

\paragraph{Assumptions.}
\begin{itemize}
  \item A square block in index space is used; with $\Delta v=0.05$\,\unit{m/s}\assumed\ and
    $\Delta\theta=0.1^\circ$\assumed\ this fixes the aspect ratio of the tolerance box.
  \item Cleanliness is checked only at the cell centres of the block.
  \item One aerodynamic set, one tilt and one $(h,d)$ at a time.
\end{itemize}

% ---------------------------------------------------------------------------------------------
\subsection{Many conditions: the robust intersection and the field-tuned hood}
\label{sec:robust_spec-hood}

\paragraph{Motivation.}
The square above is for the nominal coefficients.
On the field the drag and lift of a real POLLEN, and the HIVE tilt, can be anywhere inside their
bounds.
A fixed hood fixes the elevation angle; the wheel speed is the one thing the team can tune at the
field.
This subsection asks which hood angle leaves a usable speed under every condition.

\paragraph{Conditions.}
A condition $c$ is one aerodynamic variant combined with one HIVE tilt.
\code{shooter\_spec.variants} gives the nominal set and the drag bounds (Cd\_low, Cd\_high); for
$S\ne0$ it adds the lift bounds (Cl\_low, Cl\_high) and fast spin decay (k\_high).
The tilts are $28.80^\circ$\derived, $30^\circ$\confirmed\ and $31.77^\circ$\derived\ (the field
acceptance lip-height tolerance, solved exactly for the tilt).
That gives $3\times3=9$\derived\ conditions for $S=0$ and $6\times3=18$\derived\ for $S\ne0$.
Each condition has its own mask $M_c(v_i,\theta_j)$.

\paragraph{Robust intersection.}
The strictest requirement is a single $(v,\theta)$ that is clean in every condition:
\begin{equation}
  M_{\mathrm{rob}}(v_i,\theta_j) = \prod_{c} M_c(v_i,\theta_j).
  \label{eq:robust_spec-robust}
\end{equation}
On the unfavourable side this mask is empty in every one of the 600\predicted\ rows of
\code{data/spec\_robust.csv}: no single speed and angle works for all drag, lift and tilt bounds at
once.
This is why the specification must rely on field tuning.

\paragraph{Elevation band.}
A real hood angle is set to within $\pm\tau$ (\unit{deg}).
With $n_b=\mathrm{round}(\tau/\Delta\theta)$ (cells), a speed cell is usable at hood index $j$ only if
it is clean at every elevation in the band:
\begin{equation}
  G_c(v_i,j) = \prod_{|j'-j|\le n_b} M_c(v_i,\theta_{j'}),\qquad n_b\le j< N_\theta-n_b .
  \label{eq:robust_spec-band}
\end{equation}

\paragraph{Longest run.}
For fixed $j$, \code{longest\_run} scans $G_c(\cdot,j)$ along the speed axis.
With the run length $L_i = (L_{i-1}+1)\,G_c(v_i,j)$, $L_{-1}=0$, the longest run and its centre (a
speed index, possibly a half integer) are
\begin{equation}
  n_c(j) = \max_i L_i,\qquad
  \hat\imath_c(j) = \arg\max_i L_i - \frac{n_c(j)-1}{2}.
  \label{eq:robust_spec-run}
\end{equation}
The nominal-only fixed-hood value in \code{describe} uses only the nominal condition and picks the $j$
with the longest run (ties broken by the larger worst-case clearance), with set speed
$(v_{\mathrm{lo}}+v_{\mathrm{hi}})/2$ and tolerance $(v_{\mathrm{hi}}-v_{\mathrm{lo}})/2$ over that run.

\paragraph{Field-tuned hood.}
\code{tunable\_rows} lets each condition use its own speed.
The hood is feasible at index $j$ if every condition has at least one usable speed run.
The guaranteed run length is the worst one,
\begin{equation}
  w(j) = \min_c n_c(j),\qquad \text{feasible}\iff w(j)>0 ,
  \label{eq:robust_spec-wmin}
\end{equation}
and the spread of the run centres measures how far the team may have to re-tune,
$\Sigma(j) = \max_c \hat\imath_c(j) - \min_c \hat\imath_c(j)$ (cells; $\Delta v\,\Sigma(j)$ in \unit{m/s}).
The code ranks hood indices by a lexicographic score (run length first, small spread second; the
spread is at most $N_v-1=110$ cells, so the factor 1000 keeps the run length dominant):
\begin{equation}
  j^\ast = \arg\max_j \Big[1000\,w(j) - \Sigma(j)\Big].
  \label{eq:robust_spec-score}
\end{equation}
The reported tuning range and the speed tolerance that every condition keeps at the chosen hood are
\begin{equation}
  v_{\mathrm{tune}}\in\Big[v_{\min}+\Delta v\min_c\hat\imath_c(j^\ast),\
                          v_{\min}+\Delta v\max_c\hat\imath_c(j^\ast)\Big],\qquad
  \delta_{v,\mathrm{tune}} = \frac{w(j^\ast)-1}{2}\,\Delta v .
  \label{eq:robust_spec-tune}
\end{equation}
Note that the masks are not intersected here: Eq.~\eqref{eq:robust_spec-wmin} needs a run to exist in each
condition separately, and the runs may sit at different speeds.
% Code/doc difference: the task wording speaks of masks "intersected over conditions". The code
% (tunable_rows) takes min_c of run lengths with per-condition speeds; the true intersection is
% spec_rows / spec_robust.csv, and it is empty. This text follows the code.

\paragraph{Worked example.}
At $h=0.737$\,\unit{m}\assumed\ and $d=0.75$\,\unit{m}\assumed\ with $\tau=0.2^\circ$\assumed\
($n_b=2$\derived):
\begin{itemize}
  \item No spin: the hood is $\theta_{\mathrm{hood}}=80.9^\circ$\predicted.
    The nine run centres lie between 6.7\,\unit{m/s}\predicted\ (Cd\_low) and
    7.4\,\unit{m/s}\predicted\ (Cd\_high); the nominal-drag centres are
    6.875--6.9\,\unit{m/s}\predicted.
    The worst run is $w=1$\predicted\ cell, so $\delta_{v,\mathrm{tune}}=0$\derived: each condition has
    exactly one tuned speed that works.
  \item Backspin $S\in\{0.5,1.0,1.5\}$: no hood index is feasible at $\tau=0.2^\circ$\assumed.
    With $\tau=0.1^\circ$\assumed\ ($n_b=1$\derived) the code picks $80.4^\circ$\predicted\ for both
    $S=0.5$ and $S=1.0$, with $w=3$\predicted, i.e.\
    $\delta_{v,\mathrm{tune}}=0.05$\,\unit{m/s}\derived.
    The no-spin twin at the same band picks $84.0^\circ$\predicted\ with $w=4$\predicted\
    ($\delta_{v,\mathrm{tune}}=0.075$\,\unit{m/s}\derived).
    Neither search applies the $8.909$\,\unit{m/s}\predicted\ grip limit of
    Section~\ref{sec:launcher:grip}, and both tuning ranges reach above it (to 9.125\,\unit{m/s}\predicted\
    with backspin, 9.575\,\unit{m/s}\predicted\ without).
    With the speeds capped there, backspin keeps $w=2$\predicted\ at $80.3^\circ$\predicted\ and the
    twin $w=3$\predicted\ at $82.0^\circ$\predicted.
  \item With $\tau=0.5^\circ$\assumed\ no row of the sweep is feasible for any spin.
\end{itemize}
Without the band ($\tau=0$), every condition has clean speeds at the backspin hood $80.4^\circ$\predicted\
(conditions of $S=0.5$ and $S=1$ together);
the lowest clean speed per condition ranges from 6.85\,\unit{m/s}\predicted\ to 8.6\,\unit{m/s}\predicted\
(the high-lift set needs the most speed).
For no spin at $80.9^\circ$\predicted\ the union of the nine clean speed sets is
6.6--7.0\,\unit{m/s}\predicted\ and 7.2--7.5\,\unit{m/s}\predicted.
An exit speed range of about 6.5--9.0\,\unit{m/s}\predicted\ therefore contains a clean speed for
every condition at the chosen hood.
The nominal-only fixed hood of \code{describe} at $\tau=0.2^\circ$\assumed\ is $82.8^\circ$\predicted\
for no spin (7.75--7.8\,\unit{m/s}\predicted) and $80.9^\circ$\predicted\ for $S=0.5$
(7.6--7.65\,\unit{m/s}\predicted) and $S=1$ (7.8--7.85\,\unit{m/s}\predicted); at
$\tau=0.5^\circ$\assumed\ it finds none.
At $80.9^\circ$\predicted\ the high-lift conditions of $S=0.5$ and $S=1$ (all three tilts) have no
clean speed even at $\tau=0$, so this nominal-only hood does not carry over to the high lift set.
% Code/doc difference: docs/08_shooter_spec.md says a fixed hood at +-0.2 deg shrinks the window to
% a single 0.05 m/s cell. In the code this holds for S = 0 only; for backspin tunable_rows finds no
% feasible hood at +-0.2 deg (only at +-0.1 deg).
% Code/doc difference: docs/08 "about 81 deg (no spin)" matches the tunable hood 80.9 deg; the
% nominal-only hood02 value for S = 0 is 82.8 deg.

\paragraph{Assumptions.}
\begin{itemize}
  \item The team can re-tune wheel speed at the field but not the hood angle.
  \item The conditions are one-at-a-time changes from the nominal set (drag low/high; for $S\ne0$ also
    lift low/high and fast spin decay) at three tilts, not joint corners and not a measured distribution.
    The floor lift set of Table~\ref{tab:flight-clsets} ($c_1=0.02$\assumed, $C_{L,\max}=0.04$\assumed),
    which \code{research/params\_lit.json} recommends as the pessimistic set, is not included.
  \item Topspin ($S=-1$\assumed) is feasible at $\tau=0.2^\circ$\assumed\ in the sweep
    (for example $82.0^\circ$\predicted, 6.5--7.2\,\unit{m/s}\predicted) but is not the chosen
    configuration.
\end{itemize}

% ---------------------------------------------------------------------------------------------
\subsection{Release errors as a Gaussian blur}
\label{sec:robust_spec-smooth}

\paragraph{Motivation.}
The square and the band treat the tolerance as a hard box.
Real release errors are random, so the right question is the probability of a clean shot when the
shooter aims at a cell.
For independent Gaussian errors this probability is the clean mask blurred by a Gaussian.

\paragraph{Expectation as convolution.}
Let the actual distance, speed and elevation be $d+\delta_d$, $v(1+\varepsilon_v)$ and
$\theta+\delta_\theta$, with $\delta_d\sim\mathcal N(0,\sigma_d^2)$ (\unit{m}),
$\varepsilon_v\sim\mathcal N(0,\sigma_{v,\mathrm{rel}}^2)$ (dimensionless) and
$\delta_\theta\sim\mathcal N(0,\sigma_\theta^2)$ (\unit{deg}), all independent.
The probability of a clean shot is
\begin{equation}
  p(d,v,\theta) = \mathbb E\big[M(d+\delta_d,\ v(1+\varepsilon_v),\ \theta+\delta_\theta)\big].
  \label{eq:robust_spec-expect}
\end{equation}
Near a reference speed $v_{\mathrm{ref}}$ (\unit{m/s}) the speed error is
$v\varepsilon_v\approx v_{\mathrm{ref}}\varepsilon_v$, an additive error of standard deviation
$\sigma_v = \sigma_{v,\mathrm{rel}}v_{\mathrm{ref}}$ (\unit{m/s}).
Writing the expectation as an integral against the product density $\phi$ gives a convolution:
\begin{equation}
  \begin{split}
  p(d,v,\theta) &= \iiint M(d',v',\theta')\,
  \phi_{\sigma_d}(d'-d)\,\phi_{\sigma_v}(v'-v)\,\phi_{\sigma_\theta}(\theta'-\theta)\,
  \mathrm dd'\,\mathrm dv'\,\mathrm d\theta'\\
  &= (M*\phi)(d,v,\theta).
  \end{split}
  \label{eq:robust_spec-conv}
\end{equation}
On the grid, \code{smooth\_prob} evaluates this with \code{scipy.ndimage.gaussian\_filter}, whose
widths are in cells:
\begin{equation}
  s_d = \frac{\sigma_d}{\Delta d},\qquad s_v = \frac{\sigma_{v,\mathrm{rel}}\,v_{\mathrm{ref}}}{\Delta v},\qquad
  s_\theta = \frac{\sigma_\theta}{\Delta\theta}.
  \label{eq:robust_spec-sigcells}
\end{equation}
Cells outside the grid count as dirty (\code{mode="constant"}, value 0) and the kernel is cut at
$3\sigma$.
The speed width is applied piecewise: the mask is filtered three times and each speed row keeps the
result for its nearest reference speed,
\begin{equation}
  v_{\mathrm{ref}}(v_i) = \arg\min_{r\in\{3.5,\,6.0,\,8.5\}} |v_i-r| .
  \label{eq:robust_spec-vref}
\end{equation}
The reference speeds 3.5, 6.0 and 8.5\,\unit{m/s}\assumed\ are code constants.
The best aim at each distance, for a launcher that cannot exceed $v_{\mathrm{cap}}$ (\unit{m/s}), is
\begin{equation}
  P_{\mathrm{kin}}(d) = \max_{v_i\le v_{\mathrm{cap}},\ \theta_j} p(d,v_i,\theta_j).
  \label{eq:robust_spec-pkin}
\end{equation}

\paragraph{Error sets.}
The release error sets of \code{ERR\_SETS} are design assumptions, not measurements.
Listed as $\sigma_{v,\mathrm{rel}}$, $\sigma_\theta$, $\sigma_d$, they are:
tight 0.005\assumed, $0.15^\circ$\assumed, 0.010\,\unit{m}\assumed;
nominal 0.015\assumed, $0.40^\circ$\assumed, 0.025\,\unit{m}\assumed;
loose 0.030\assumed, $0.80^\circ$\assumed, 0.050\,\unit{m}\assumed.
The azimuth and lateral spreads $\sigma_\psi$, $\sigma_y$ are $0.25^\circ$\assumed, 0.010\,\unit{m}\assumed\
(tight), $0.50^\circ$\assumed, 0.025\,\unit{m}\assumed\ (nominal) and $1.00^\circ$\assumed,
0.050\,\unit{m}\assumed\ (loose).

\paragraph{Worked example.}
On the spec grid ($\Delta d = 0.05$\,\unit{m}\assumed) with the nominal errors,
$s_d = 0.025/0.05 = 0.5$\derived\ cells, $s_\theta = 0.40/0.1 = 4.0$\derived\ cells and
$s_v = 0.015\times6.0/0.05 = 1.8$\derived\ cells near 6\,\unit{m/s} (1.05\derived\ near 3.5 and
2.55\derived\ near 8.5\,\unit{m/s}).
At $h=0.737$\,\unit{m}\assumed\ the best smoothed probability with no spin is
0.424\predicted\ at $d=0.75$\,\unit{m}\assumed\ (aim 7.2\,\unit{m/s}\predicted, $81.6^\circ$\predicted),
0.430\predicted\ at $d=0.80$\,\unit{m}\assumed\ (7.05\,\unit{m/s}\predicted, $80.9^\circ$\predicted) and
0.362\predicted\ at $d=1.15$\,\unit{m}\assumed\ (aim 9.1\,\unit{m/s}\predicted, above the grip limit;
0.356\predicted\ with $v_{\mathrm{cap}}=8.909$\,\unit{m/s}\predicted\ of Section~\ref{sec:launcher:grip}).
With $S=1$ it is 0.378\predicted\ at $d=0.75$\,\unit{m}\assumed\ (8.2\,\unit{m/s}\predicted,
$81.3^\circ$\predicted) and 0.412\predicted\ at $d=0.80$\,\unit{m}\assumed\ (7.25\,\unit{m/s}\predicted,
$79.7^\circ$\predicted), below no spin at both distances.
Two points stand out.
First, the smoothed optimum sits near $80$--$82^\circ$, not at the $83.9^\circ$ no-spin square of
Section~\ref{sec:robust_spec-square}: a blur of $4$ cells in elevation rewards the part of the window that is
deep in elevation, while the square rewards the part that is wide in both directions.
Second, even the best aim is clean only about four times in ten with nominal release errors.
On the robustness grid of \code{robustness.py} ($\Delta d=0.01$\,\unit{m}\assumed) the same errors give
$s_d=2.5$\derived\ cells.

\paragraph{Assumptions.}
\begin{itemize}
  \item Errors in $d$, $v$ and $\theta$ are Gaussian, independent and the same for every shot.
  \item The relative speed error is approximated by three fixed additive widths.
  \item The mask is for one aerodynamic condition; physical spread is added in Section~\ref{sec:robust_spec-mc}.
\end{itemize}

% ---------------------------------------------------------------------------------------------
\subsection{Azimuth: why the in-plane distance scales with $\cos\psi$}
\label{sec:robust_spec-azimuth}

\paragraph{Motivation.}
A robot that is not square to the HIVE shoots at an azimuth $\psi$ (\unit{deg} or \unit{rad}) away
from the $x$ axis.
Re-running the 3D solver for every azimuth would be expensive.
\code{clean\_masks} instead reuses the 2D paths and scales their horizontal coordinate and velocity
by $\cos\psi$.
This subsection shows why that is exact in still air.

\paragraph{Rotation about the vertical.}
Let
\begin{equation}
  \mathsf R_\psi = \begin{pmatrix}\cos\psi & -\sin\psi & 0\\ \sin\psi & \cos\psi & 0\\ 0&0&1\end{pmatrix}
  \label{eq:robust_spec-Rz}
\end{equation}
be the rotation by $\psi$ about $\hat{\vect z}$.
It preserves lengths, $|\mathsf R_\psi\vect a| = |\vect a|$, it commutes with the cross product,
$(\mathsf R_\psi\vect a)\times(\mathsf R_\psi\vect b) = \mathsf R_\psi(\vect a\times\vect b)$ (true for
every proper rotation), and it leaves the vertical unchanged, $\mathsf R_\psi\hat{\vect z}=\hat{\vect z}$.

\paragraph{The acceleration is equivariant.}
In still air the 3D acceleration of \code{integrate3d} is
\begin{equation}
  \vect a(\vect v,\vect\omega) = -q\,C_D\,|\vect v|\,\vect v
  + q\,C_L(S)\,|\vect v|^2\,\frac{\vect\omega\times\vect v}{|\vect\omega\times\vect v|} - g\,\hat{\vect z},
  \qquad q=\frac{\rho\pi R^2}{2m},
  \label{eq:robust_spec-acc3d}
\end{equation}
with air density $\rho$ (\unit{kg/m\textsuperscript{3}}), ball mass $m$ (\unit{kg}), drag and lift coefficients
$C_D$, $C_L$ (dimensionless), $q$ in \unit{1/m} and $g$ in \unit{m/s\textsuperscript{2}}.
(The code regularises the two denominators with $\varepsilon=10^{-6}$\assumed; this does not affect the
argument.)
The spin parameter $S=|\vect\omega|R/|\vect v|$ is a ratio of lengths, so it is unchanged by the
rotation.
Applying the three properties term by term:
\begin{equation}
  \vect a(\mathsf R_\psi\vect v,\mathsf R_\psi\vect\omega)
  = -qC_D|\vect v|\,\mathsf R_\psi\vect v
  + qC_L|\vect v|^2\,\frac{\mathsf R_\psi(\vect\omega\times\vect v)}{|\vect\omega\times\vect v|}
  - g\,\mathsf R_\psi\hat{\vect z}
  = \mathsf R_\psi\,\vect a(\vect v,\vect\omega).
  \label{eq:robust_spec-equiv}
\end{equation}
The spin decays as $\vect\omega(t)=\vect\omega_0e^{-kt}$ about a fixed axis, with decay rate
$k$ (\unit{1/s}), so the rotated spin stays the rotated spin at all times.

\paragraph{The backspin axis rotates with the shot.}
The launch direction is $\hat{\vect v}_0 = (\cos\theta\cos\psi,\ \cos\theta\sin\psi,\ \sin\theta)$ and
\code{backspin\_axis} takes the unit vector along $\hat{\vect v}_0\times\hat{\vect z}$:
\begin{equation}
  \begin{split}
  \hat{\vect v}_0\times\hat{\vect z} &= (\cos\theta\sin\psi,\ -\cos\theta\cos\psi,\ 0),\\
  \hat{\vect e}_\omega &= (\sin\psi,\ -\cos\psi,\ 0) = \mathsf R_\psi(0,-1,0)^{\mathsf T}.
  \end{split}
  \label{eq:robust_spec-axis}
\end{equation}
At $\psi=0$ this is $-\hat{\vect y}$, and $(-\hat{\vect y})\times(v_x,0,v_z)=(-v_z,0,v_x)$, which is
the 2D Magnus direction of Section~\ref{sec:flight}.

\paragraph{The path rotates rigidly.}
Let $\vect r_2(t)=(x(t),0,z(t))$ be the 2D path launched at $\psi=0$.
By Eqs.~\eqref{eq:robust_spec-equiv}--\eqref{eq:robust_spec-axis}, $\mathsf R_\psi\vect r_2(t)$ satisfies the same
equation of motion as the 3D path launched at $\psi$, from the rotated initial state.
Solutions of an ODE with a Lipschitz right-hand side are unique, so
\begin{equation}
  \vect r_3(t) = \mathsf R_\psi\,\vect r_2(t) = \big(x(t)\cos\psi,\ x(t)\sin\psi,\ z(t)\big).
  \label{eq:robust_spec-proj}
\end{equation}
The HIVE cross-section does not change along the pivot axis $y$ (inside the CELL width), so only the
$x$--$z$ projection matters for the clean test.
That projection keeps $z$ and scales the horizontal coordinate and velocity:
\begin{equation}
  X(t) = x(t)\cos\psi,\qquad V_X(t) = v_x(t)\cos\psi,\qquad Z(t)=z(t),\qquad V_Z(t)=v_z(t),
  \label{eq:robust_spec-projX}
\end{equation}
which is exactly what \code{clean\_masks} applies before the gate test.
Wind breaks the argument: with $\vect v_r = \vect v-\vect w$ and a fixed wind $\vect w$ (\unit{m/s}),
$\mathsf R_\psi\vect w\ne\vect w$ in general, so the 3D Monte Carlo integrates with wind directly.

\paragraph{Azimuth from the lateral offset.}
A shot aimed at lateral offset $y$ (\unit{m}) from the CELL centre reaches the gate plane at
horizontal distance $d\pm x_g$, where $x_g=0.401$\,\unit{m}\derived\ is the gate position in the HIVE
frame.
The robustness maps use
\begin{equation}
  \psi = \arctan\frac{|y|}{d\pm x_g},\qquad
  s_{\mathrm{eff}} = \min(0.3\tan\psi,\ 0.2)\ \unit{m},
  \label{eq:robust_spec-psi}
\end{equation}
with azimuth bins $0^\circ$ to $60^\circ$\assumed\ in $10^\circ$\assumed\ steps; $s_{\mathrm{eff}}$ is an
assumed effective sideways shift of the mouth seen at an angle (coefficients 0.3\assumed\ and
0.2\,\unit{m}\assumed).

\paragraph{Lateral factor.}
The sideways miss at the gate is $e_y$ (\unit{m}).
An azimuth error $\delta_\psi\sim\mathcal N(0,\sigma_\psi^2)$ (\unit{rad}) over a lever arm
$L=\sqrt{(d\pm x_g)^2+y^2}$ (\unit{m}, the horizontal release-to-gate distance used by
\code{position\_prob}; $L=d+x_g$ for an unfavourable shot with $y=0$) moves the ball by
$L\tan\delta_\psi\approx L\delta_\psi$.
Adding a lateral release error $\sigma_y$ (\unit{m}) and a wind drift $\sigma_{\mathrm{dr}}$ (\unit{m}),
all independent and Gaussian, the variances add:
\begin{equation}
  \sigma_{\mathrm{lat}}^2 = (L\sigma_\psi)^2 + \sigma_y^2 + \sigma_{\mathrm{dr}}^2 .
  \label{eq:robust_spec-latsig}
\end{equation}
The ball fits if $|e_y|<w_{\mathrm{lat}}$, with half-width $w_{\mathrm{lat}}$ (\unit{m}).
For $e_y\sim\mathcal N(0,\sigma_{\mathrm{lat}}^2)$, substituting $u=e/(\sqrt2\sigma_{\mathrm{lat}})$,
\begin{equation}
  \begin{split}
  \Pr(|e_y|<w_{\mathrm{lat}})
  &= \int_{-w_{\mathrm{lat}}}^{w_{\mathrm{lat}}}\frac{e^{-e^2/(2\sigma_{\mathrm{lat}}^2)}}{\sqrt{2\pi}\,\sigma_{\mathrm{lat}}}\,\mathrm de
  = \frac{2}{\sqrt\pi}\int_0^{w_{\mathrm{lat}}/(\sqrt2\sigma_{\mathrm{lat}})}e^{-u^2}\mathrm du\\
  &= \operatorname{erf}\!\Big(\frac{w_{\mathrm{lat}}}{\sqrt2\,\sigma_{\mathrm{lat}}}\Big).
  \end{split}
  \label{eq:robust_spec-erf}
\end{equation}
\code{position\_prob} combines it with the kinematic probability, treating the two as independent:
\begin{equation}
  P_{\mathrm{pos}}(d,\psi) = P_{\mathrm{kin}}(d,\psi)\,\operatorname{erf}\!\Big(\frac{w_{\mathrm{lat}}}{\sqrt2\,\sigma_{\mathrm{lat}}}\Big).
  \label{eq:robust_spec-ppos}
\end{equation}
The code uses $w_{\mathrm{lat}} = 0.10$\,\unit{m}\assumed\ (low), 0.15\,\unit{m}\assumed\ (primary) and
0.2185\,\unit{m}\derived\ (high, the CELL half-width 0.254\,\unit{m}\confirmed\ minus $R$).

\paragraph{Worked example.}
For the 3D path with params\_lit nominal aerodynamics, $v=7.25$\,\unit{m/s}\assumed,
$\theta=79.7^\circ$\assumed, $S=1$\assumed\ and $\psi=30^\circ$\assumed, the backspin axis is
$(0.5,-0.866,0)$\derived.
Comparing \code{integrate3d} with the 2D solution rotated by Eq.~\eqref{eq:robust_spec-proj}, the largest
differences are $1.1\times10^{-7}$\,\unit{m}\predicted\ in $x$, $1.1\times10^{-16}$\,\unit{m}\predicted\ in
$y$ and 0\predicted\ in $z$: the projection is exact to round-off.
On this path the apex is 2.339\,\unit{m}\predicted\ above the release point; there $v_z=0$ but the ball
still moves horizontally at $v_x=0.903$\,\unit{m/s}\predicted, so it is not at rest at the apex.

For the lateral factor take $d=0.81$\,\unit{m}\assumed, so $L=0.81+0.401=1.211$\,\unit{m}\derived, and the
nominal $\sigma_\psi=0.5^\circ=0.00873$\,\unit{rad}\assumed\ and $\sigma_y=0.025$\,\unit{m}\assumed.
Then $L\sigma_\psi=0.0106$\,\unit{m}\derived\ and, without wind,
$\sigma_{\mathrm{lat}}=\sqrt{0.0106^2+0.025^2}=0.0271$\,\unit{m}\derived; the factor is
1.000\derived\ for $w_{\mathrm{lat}}=0.15$\,\unit{m}\assumed.
With the wind drift $\sigma_{\mathrm{dr}}=0.0615$\,\unit{m}\derived\ from \code{wind\_drift\_sd}
(module-default coefficients \code{data/aero\_params.json}; 0.0634\,\unit{m}\derived\ with
\code{research/params\_lit.json}),
$\sigma_{\mathrm{lat}}=0.0672$\,\unit{m}\derived\ and the factor drops to
$\operatorname{erf}(0.15/(\sqrt2\cdot0.0672))=0.974$\derived, or 0.863\derived\ for
$w_{\mathrm{lat}}=0.10$\,\unit{m}\assumed.
Sideways aim is a small loss next to the in-plane window.

\paragraph{Assumptions.}
\begin{itemize}
  \item Still air for the projection; wind enters only through $\sigma_{\mathrm{dr}}$ and the 3D Monte Carlo.
  \item The HIVE section is uniform along the pivot axis across the CELL width.
  \item Lateral and in-plane outcomes are independent in Eq.~\eqref{eq:robust_spec-ppos}.
  \item Azimuth bins stop at $60^\circ$\assumed\ (\code{PSI\_MAX}).
\end{itemize}

% ---------------------------------------------------------------------------------------------
\subsection{Three-dimensional Monte Carlo}
\label{sec:robust_spec-mc}

\paragraph{Motivation.}
The blur of Section~\ref{sec:robust_spec-smooth} covers release errors in one aerodynamic condition.
A real match mixes everything at once: release errors, ball-to-ball spread, unknown coefficients,
air movement and the HIVE pose.
\code{mc\_shots} draws all of these per shot and flies each shot in 3D.

\paragraph{Per-shot sampling.}
Each shot $n$ draws independent standard normals $\varepsilon_{1},\varepsilon_2,\dots\sim\mathcal N(0,1)$.
The release state is
\begin{align}
  v &= v_{\mathrm{aim}}(1+\sigma_{v,\mathrm{rel}}\varepsilon_1)(1+b) - \frac{\partial v}{\partial n_w}\,\Delta n_w\,f,
  \label{eq:robust_spec-mc-v}\\
  \theta &= \theta_{\mathrm{aim}}+\sigma_\theta\varepsilon_2,\qquad \psi = \sigma_\psi\varepsilon_3,
  \label{eq:robust_spec-mc-ang}\\
  \vect r_0 &= \big(-(d+\sigma_d\varepsilon_4),\ \sigma_y\varepsilon_5,\ h+\sigma_h\varepsilon_6\big),
  \label{eq:robust_spec-mc-pos}
\end{align}
where $b$ (dimensionless) is a uniform speed bias used by the battery-sag scenarios,
$\partial v/\partial n_w$ (\unit{m/s} per rpm) is the launcher-curve slope, $\Delta n_w$ (rpm) the
wheel-speed drop after a shot and $f$ a uniform recovery-deficit fraction for the rapid-fire
scenarios; in the nominal scenario $b=f=0$.
The release-height error is $\sigma_h = 0.005$\,\unit{m}\assumed.
The ball and the air are drawn as
\begin{equation}
  \begin{aligned}
    m &= \operatorname{clip}(\bar m+\sigma_m\varepsilon_7,\ m_{\mathrm{lo}},m_{\mathrm{hi}}), &
    R &= \operatorname{clip}(\bar R+\sigma_R\varepsilon_8,\ R_{\mathrm{lo}},R_{\mathrm{hi}}),\\
    C_D &= \bar C_D(1+0.05\,\varepsilon_9), &
    (c_1,C_{L,\max}) &= (\bar c_1,\bar C_{L,\max})(1+0.10\,\varepsilon_{10}),\\
    \vect w &= \sigma_w(\varepsilon_{11},\varepsilon_{12},0), &
    S &\sim\mathcal U(S_{\mathrm{lo}},S_{\mathrm{hi}}),
  \end{aligned}
  \label{eq:robust_spec-mc-phys}
\end{equation}
with $\sigma_m=0.5$\,\unit{g}\assumed, $\sigma_R=0.25$\,\unit{mm}\assumed, the coefficient spreads
5\,\%\assumed\ and 10\,\%\assumed\ (one factor shared by $c_1$ and $C_{L,\max}$, so $S_{\mathrm{sat}}$
does not move), a wind of $\sigma_w=0.30$\,\unit{m/s}\assumed\ per horizontal component held constant
during the shot, and a spin parameter uniform on $[1.03, 1.27]$\predicted\ (the contact model's hood
band); the release-error-only scenario fixes $S=1.0597$\predicted.
The HIVE is drawn per shot with tilt $\beta=\operatorname{clip}(30^\circ+0.75^\circ\varepsilon_{13},\
28.80^\circ, 31.77^\circ)$ (spread $0.75^\circ$\assumed, clipped to the derived tilt bounds), pivot height
$1.11633$\,\unit{m}\confirmed\ with spread 0.005\,\unit{m}\assumed\ clipped to
1.09093--1.14173\,\unit{m}\confirmed\ (the manual's $\pm1$\,in field tolerance), and an $x$ offset of
spread 0.01\,\unit{m}\assumed.
The initial spin is backspin about the rotated axis of Eq.~\eqref{eq:robust_spec-axis}:
\begin{equation}
  \vect\omega_0 = \hat{\vect e}_\omega(\psi)\,\frac{S\,v}{R}.
  \label{eq:robust_spec-mc-omega}
\end{equation}
The path is integrated with Eq.~\eqref{eq:robust_spec-acc3d} using the relative velocity
$\vect v_r=\vect v-\vect w$.

\paragraph{Scoring one shot.}
Where the path (with $x$ measured from the drawn HIVE position) last crosses the gate plane
$x=x_g$ between samples $k$ and $k+1$, the lateral position is interpolated linearly:
\begin{equation}
  y_g = y_k + (y_{k+1}-y_k)\,\frac{x_g-x_k}{x_{k+1}-x_k}.
  \label{eq:robust_spec-yg}
\end{equation}
If $|y_g| > w_{\mathrm{cell}}-R$, with CELL half-width $w_{\mathrm{cell}}=0.254$\,\unit{m}\confirmed,
the shot hits a side wall (or misses when it never crosses the plane).
Otherwise a HIVE with the drawn tilt, pivot height and ball radius, offset sideways by $|y_g|$, runs
\code{classify\_path} on the in-plane path.
The shot counts as clean only if the outcome is \code{clean}.
Out of $n$ shots, $k$ are clean, and the point estimate is
\begin{equation}
  \hat p = \frac{k}{n}.
  \label{eq:robust_spec-phat}
\end{equation}

\paragraph{Worked example.}
The stored run \code{data/robust\_mc.csv} aims at $d=0.81$\,\unit{m}\predicted,
$v=7.25$\,\unit{m/s}\predicted, $\theta=79.7^\circ$\predicted\ with backspin and $n=1500$\assumed\ shots
per scenario.
Its nominal scenario has $k=455$\predicted, $\hat p=0.303$\predicted, Wilson interval
$[0.281, 0.327]$\derived, and a gate-lateral spread of 0.049\,\unit{m}\predicted.
The other scenarios (all at $n=1500$\assumed):
\begin{itemize}
  \item release errors only (fixed physics, no wind): $\hat p=0.362$\predicted;
  \item tight release errors: 0.411\predicted; loose: 0.157\predicted;
  \item wind $\sigma_w=1.0$\,\unit{m/s}\assumed: 0.111\predicted;
  \item true lift at the low or high bound while the aim was tuned for nominal lift:
    0.097\predicted\ and 0.054\predicted; re-tuned for low lift ($\theta=80.7^\circ$\predicted):
    0.325\predicted;
  \item the no-spin twin aimed at $81.4^\circ$\predicted: 0.317\predicted\ $[0.294, 0.341]$\derived;
  \item the favourable side for comparison: 0.996\predicted.
\end{itemize}
The spin and no-spin intervals overlap, so at this aim backspin gives no measurable gain.
Using the wrong lift coefficient costs far more than the release scatter, which is why the
specification relies on tuning at the field.
A fresh run of 400\assumed\ shots with seed 1000\assumed\ and the same code gives
$k=110$\predicted, $\hat p=0.275$\predicted, Wilson $[0.234, 0.321]$\derived; it overlaps the stored
result.
The two most common failures there are the outer roof (147\predicted\ shots) and the floor
(62\predicted).
% Code/doc difference: docs/06_required_vs_achievable.md and the robust_mc / robust_map CSVs were
% produced with data/aero_params.json (older first-pass set). The robustness CLI default is now
% research/params_lit.json, but data/params_lit/ is empty: no params_lit rerun exists. The numbers
% above are the aero_params.json results.
% Code/doc difference: flight.AERO_PARAMS_PATH (module level) still defaults to
% data/aero_params.json unless PM001_AERO_PATH is set, so code that imports robustness without the
% CLI gets the old set. The 400-shot rerun above used that module default to match the CSV.
% Code/doc difference: docs/06 says the mass is clipped to 23.5-26.5 g (aero_params.json); the
% params_lit upper bound is 27.5 g.
These Monte Carlo results were computed with the older first-pass coefficient file
\code{data/aero\_params.json}; a rerun with \code{research/params\_lit.json} has not been made.
The stored file also predates the current scenario list of \code{mc\_scenarios}: it has no no-spin
shots at the backspin aim, and (as \code{data/aero\_params.json} has no floor set) no rows for the
floor lift set.
% The clip above is the current code (fz.TILT_BOUNDS 28.80-31.77); squashed git history cannot show
% which clip produced data/robust_mc.csv.

\paragraph{Assumptions.}
\begin{itemize}
  \item All spreads in Eqs.~\eqref{eq:robust_spec-mc-v}--\eqref{eq:robust_spec-mc-phys} are assumed, independent and
    Gaussian or uniform; none is measured on the robot.
  \item The wind is constant during one shot and horizontal.
  \item The aerodynamic coefficients are assumed literature-based ranges, not proven bounds and not
    POLLEN measurements.
  \item Shots are independent trials, which Section~\ref{sec:robust_spec-wilson} needs.
\end{itemize}

% ---------------------------------------------------------------------------------------------
\subsection{How sure is a success rate? The Wilson interval}
\label{sec:robust_spec-wilson}

\paragraph{Motivation.}
A rate of $k$ clean shots out of $n$ is itself random.
The simple interval $\hat p\pm z\sqrt{\hat p(1-\hat p)/n}$ fails badly for small $n$ or for
$\hat p$ near 0 or 1: for one clean shot out of one it gives $[1,1]$.
\code{wilson} uses the score interval instead, which is the right tool for field tests.

\paragraph{Derivation.}
For $n$ independent shots with true probability $p$, $k$ is binomial, and for moderate $n$
\begin{equation}
  \frac{\hat p-p}{\sqrt{p(1-p)/n}}\approx\mathcal N(0,1).
  \label{eq:robust_spec-wilson-z}
\end{equation}
The Wilson interval is the set of $p$ for which this score is within $\pm z$, where
$z=\Phi^{-1}(0.975)=1.959964$\derived\ for 95\,\%.
Note that the variance uses the unknown $p$, not $\hat p$.
The boundary satisfies
\begin{equation}
  (\hat p-p)^2 = \frac{z^2}{n}\,p(1-p).
  \label{eq:robust_spec-wilson-quad}
\end{equation}
Expanding both sides and collecting powers of $p$ gives a quadratic:
\begin{equation}
  \Big(1+\frac{z^2}{n}\Big)p^2 - \Big(2\hat p+\frac{z^2}{n}\Big)p + \hat p^2 = 0 .
  \label{eq:robust_spec-wilson-poly}
\end{equation}
Its discriminant simplifies as
\begin{equation}
  \Big(2\hat p+\frac{z^2}{n}\Big)^2 - 4\Big(1+\frac{z^2}{n}\Big)\hat p^2
  = \frac{4z^2}{n}\hat p(1-\hat p) + \frac{z^4}{n^2}
  = 4z^2\Big(\frac{\hat p(1-\hat p)}{n}+\frac{z^2}{4n^2}\Big),
  \label{eq:robust_spec-wilson-disc}
\end{equation}
and the two roots, divided top and bottom by 2, are the interval:
\begin{equation}
  p_\pm = \frac{\hat p + \dfrac{z^2}{2n} \pm z\sqrt{\dfrac{\hat p(1-\hat p)}{n}+\dfrac{z^2}{4n^2}}}{1+\dfrac{z^2}{n}} .
  \label{eq:robust_spec-wilson}
\end{equation}
The code clips both ends to $[0,1]$.
The centre $(\hat p+z^2/2n)/(1+z^2/n)$ is pulled toward $1/2$, and the width never collapses to
zero, even at $\hat p = 0$ or 1.

\paragraph{Worked example.}
For $k=455$\predicted\ of $n=1500$\assumed: $\hat p=0.3033$\derived, $z^2/n=0.00256$\derived, and
Eq.~\eqref{eq:robust_spec-wilson} gives $[0.2806, 0.3271]$\derived.
For $\hat p = 0.5$ the half-width is 0.263\derived\ at $n=10$, 0.201\derived\ at $n=20$,
0.096\derived\ at $n=100$ and 0.025\derived\ at $n=1500$.
One clean shot out of one gives $[0.207, 1]$\derived: a single shot proves only that the clean rate is
probably above about one in five, and says nothing about repeatability.
Zero clean shots out of 20 still allows a rate up to 0.161\derived.
A field test that should separate a 30\,\% shooter from a 50\,\% one therefore needs on the order of
100 shots per setting, not a handful.

\paragraph{Assumptions.}
\begin{itemize}
  \item Shots are independent with a common probability (no warm-up drift).
  \item The normal approximation in Eq.~\eqref{eq:robust_spec-wilson-z}; the Wilson form keeps good coverage
    down to small $n$.
\end{itemize}

% ---------------------------------------------------------------------------------------------
\subsection{From exit speed to wheel and ball rpm}
\label{sec:robust_spec-rpm}

\paragraph{Motivation.}
The CAD designer chooses a motor and gear ratio, so the speed requirement must be stated as a wheel
rpm, and the spin requirement as a ball rpm.

\paragraph{Wheel rpm.}
A wheel of diameter $D$ (\unit{m}) turning at $n_w$ (rpm) has surface speed
$u=\pi D\,n_w/60$ (\unit{m/s}).
The exit speed of the ball is a fraction $r=v/u$ (dimensionless) of the surface speed.
Solving for the wheel speed:
\begin{equation}
  n_w = \frac{60\,v}{r\,\pi D}.
  \label{eq:robust_spec-wheel}
\end{equation}
This is \code{wheel\_rpm}.
The ratio $r$ is not universal.
With a single powered wheel against a fixed hood the ball rolls on the hood, and the simple
rolling picture gives $r=1/2$; slip, compression and a moving upper roller change it.
The code tabulates $r\in\{0.25, 0.30, 0.42, 0.50\}$\assumed\ (grade C), and the contact model of
the previous section is the place where $r$ is predicted.

\paragraph{Ball rpm.}
From $S=\omega R/v$, the spin rate is $\omega=Sv/R$ (\unit{rad/s}); converting with
$1\,\text{rev}=2\pi$\,\unit{rad} and $1\,\text{min}=60$\,\unit{s},
\begin{equation}
  n_b = \frac{S\,v}{R}\cdot\frac{60}{2\pi}.
  \label{eq:robust_spec-spinrpm}
\end{equation}
This is \code{spin\_rpm}.
There is no point in asking for spin far above $S_{\mathrm{sat}}$: the lift coefficient saturates
and $S=1.5$ gives the same window as $S=1.0$ (Section~\ref{sec:robust_spec-grid}).

\paragraph{Worked example.}
With $D=0.098$\,\unit{m}\userval\ and $R=0.0355$\,\unit{m}\confirmed\ (the manual's ball diameter of
about 7.1\,\unit{cm}):
\begin{itemize}
  \item $v=6.5$\,\unit{m/s}\predicted:
    $n_w = 60\times6.5/(0.25\,\pi\,0.098) = 5067$\derived\ rpm at $r=0.25$\assumed, and
    4222\derived, 3016\derived, 2533\derived\ rpm at $r=0.30$, 0.42, 0.50.
  \item $v=9.0$\,\unit{m/s}\predicted: 7016\derived, 5847\derived, 4176\derived, 3508\derived\ rpm.
  \item $v=7.0$\,\unit{m/s}\predicted: 5457\derived, 4547\derived, 3248\derived, 2728\derived\ rpm.
  \item Ball spin: $S=0.5$\assumed\ at 6.5\,\unit{m/s} is
    $n_b=0.5\times6.5/0.0355\times60/(2\pi)=874$\derived\ rpm; $S=1$\assumed\ at 7.0\,\unit{m/s} is
    1883\derived\ rpm and at 9.0\,\unit{m/s} is 2421\derived\ rpm.
\end{itemize}
The full exit-speed band 6.5--9.0\,\unit{m/s} therefore asks for anything from about 2500 to 7000 rpm
at the wheel depending on $r$; measuring $r$ on the first prototype narrows this to one gear choice.

\paragraph{Assumptions.}
\begin{itemize}
  \item $r$ is constant over the speed range for a given shooter build.
  \item The wheel speed does not drop during the shot (the drop is a Monte Carlo scenario, Eq.~\eqref{eq:robust_spec-mc-v}).
\end{itemize}

% ---------------------------------------------------------------------------------------------
\subsection{The shooter specification card}
\label{sec:robust_spec-card}

\paragraph{Motivation.}
The final step collects the results above into the few numbers the CAD designer needs.
Each target traces back to an equation and a worked example of this section.
All of them are model outputs.

\begin{table}[htbp]
  \centering
  \caption{CAD-card targets for the unfavourable-side shot (single powered 98\,mm wheel, fixed hood).
    All targets are model predictions to be confirmed on the field.}
  \label{tab:robust_spec-card}
  \begin{tabular}{@{}p{0.24\linewidth}p{0.26\linewidth}p{0.42\linewidth}@{}}
    \toprule
    Quantity & Target & Where it comes from \\
    \midrule
    Release height $h$ & 0.65--0.737\,\unit{m}\predicted & Table~\ref{tab:robust_spec-nclean}; window shrinks below 0.65\,\unit{m}\predicted \\
    Distance $d$ & 0.75--0.85\,\unit{m}\predicted & no clean cell at $d\le0.70$\,\unit{m}\predicted; peak near 0.80\,\unit{m}\predicted \\
    Hood angle, no spin & about $81^\circ$\predicted & Eq.~\eqref{eq:robust_spec-score}: $80.9^\circ$\predicted\ at $\tau=0.2^\circ$\assumed \\
    Hood angle, backspin & about $80.4^\circ$\predicted & Eq.~\eqref{eq:robust_spec-score} at $\tau=0.1^\circ$\assumed, $S=0.5$\assumed\ and $1.0$\assumed\ ($80.3^\circ$\predicted\ with speeds capped at the grip limit) \\
    Exit speed range & 6.5--9.0\,\unit{m/s}\predicted & every condition has a clean speed (Sec.~\ref{sec:robust_spec-hood}) \\
    Speed repeatability & $\pm0.075$\,\unit{m/s}\predicted\ with backspin, $\pm0.05$\,\unit{m/s}\predicted\ without & Eq.~\eqref{eq:robust_spec-tol}, $k^\ast=4$\predicted\ with backspin, $k^\ast=3$\predicted\ for the no-spin twin (nominal aerodynamics) \\
    Angle repeatability & $\pm0.15^\circ$\predicted\ with backspin, $\pm0.10^\circ$\predicted\ without & Eq.~\eqref{eq:robust_spec-tol}, as for the speed \\
    Wheel speed & about 2500--7000 rpm\predicted & Eq.~\eqref{eq:robust_spec-wheel}, depends on $r$ \\
    \bottomrule
  \end{tabular}
\end{table}

\paragraph{How to read the card.}
\begin{itemize}
  \item \emph{Height and distance.}
    The window opens at $d\approx0.75$\,\unit{m}\predicted\ and is largest near 0.80\,\unit{m}\predicted;
    beyond 0.85\,\unit{m} it shrinks slowly.
    Lower release heights shrink it quickly (Table~\ref{tab:robust_spec-nclean} and the $h$ comparison in
    Section~\ref{sec:robust_spec-grid}).
  \item \emph{Hood angle.}
    The hood is chosen so that, for every drag, lift and tilt condition, some tuned wheel speed is clean
    (Eq.~\eqref{eq:robust_spec-wmin}).
    The backspin hood is slightly flatter than the no-spin hood, but backspin does not enlarge the
    window and its hood needs a tighter angle tolerance ($\pm0.1^\circ$ instead of $\pm0.2^\circ$).
  \item \emph{Exit speed.}
    The range is a tuning range, not a tolerance: the team sets one speed at the field and keeps it.
  \item \emph{Repeatability.}
    The $\pm0.075$\,\unit{m/s}\predicted\ and $\pm0.15^\circ$\predicted\ figures are the largest clean
    square of the nominal sweep (nominal coefficients, $30^\circ$\confirmed\ tilt;
    Section~\ref{sec:robust_spec-square}), and only backspin reaches them; the no-spin twin reaches
    $\pm0.05$\,\unit{m/s}\predicted\ and $\pm0.10^\circ$\predicted.
    Each square is found for one set of coefficients and one tilt.
    They are a target for shot-to-shot scatter, roughly the tight release set of the Monte Carlo
    ($\sigma_\theta=0.15^\circ$\assumed), which raises the modelled clean rate from 0.30\predicted\ to
    0.41\predicted.
    They are not a guarantee of clean entry under every condition.
\end{itemize}
% Code/doc difference: docs/08 puts the repeatability (no spin +-0.05 m/s and +-0.1 deg, backspin
% +-0.075 m/s and +-0.15 deg) on the card next to the ~81 deg hood. In the code these are the
% nominal-aero best squares at h = 0.737 m, d = 0.75 m, tilt 30 deg: S = 0 k*=3 at (8.55 m/s,
% 83.9 deg), S = 0.5 k*=4 at (9.125 m/s, 82.35 deg, above the 8.909 m/s cap), S = 1 k*=4 at
% (8.875 m/s, 81.85 deg); none is at the ~81 deg hood and none is robust over conditions. At the
% field-tuned no-spin hood (tunable_rows, tau = 0.2 deg) the guaranteed speed tolerance is 0.
% Code/doc difference: docs/08 gives the backspin hood as about 80.5-81.5 deg (band centre). The
% field-tuned hood at +-0.1 deg is 80.4 deg (80.3 deg with speeds capped at 8.909 m/s). The
% nominal-only hood (fixed_elev at +-0.2 deg) is 80.9 deg, where the Cl_high conditions have no clean
% speed. At +-0.2 deg the field-tuned search finds no backspin hood.

\paragraph{Assumptions.}
\begin{itemize}
  \item A single powered wheel with a fixed hood (an optional passive upper roller changes $r$ and
    the spin, not the geometry of the window).
  \item The coefficient ranges of \code{research/params\_lit.json} for the spec grid; the Monte Carlo
    rates were computed with \code{data/aero\_params.json}.
  \item No claim here is based on a single test shot; field confirmation needs a Wilson interval
    from many shots (Section~\ref{sec:robust_spec-wilson}).
\end{itemize}
